%% AASTeX v6.3.1 standard formatting for professional astrophysics submission
\documentclass[trackchanges]{aastex631}

\begin{document}

\title{Evidence for Empirical Kinematic Signatures of Cosmic Void Boundaries:\\
Triaxial Velocity Shear and Multi-Void Survey in the Cosmicflows-4 Catalog and DESI DR1}

\author[0000-0000-0000-0000]{Jacob J. Rasmussen}
\affiliation{Independent Researcher, Spooner, WI 54801}

\correspondingauthor{Jacob J. Rasmussen}
\email{Independent Researcher}

\begin{abstract}
We present an observational analysis of localized velocity shear at cosmic void perimeters using the Cosmicflows-4 (CF4) Extragalactic Distance Database ($N = 38,057$ total groups). To address selection bias, algorithm dependencies, and projection systematics, we evaluate a multi-void survey of $N = 10$ local cosmic structures ($N_{\text{gal}}$ ranging from 18 to 199 per boundary shell, standardized at $28.5 \le r_{\text{void}} \le 34.5\text{ Mpc}$ with a $6\text{ Mpc}$ thickness). By projecting Earth-centric 1D peculiar velocities into 3D void-centric reference frames with 500-iteration Monte Carlo bootstrap error propagation and permutation testing, we isolate radial expansion ($V_r$) and tangential shear ($V_t$). We define phase-drift ($\Phi_D$) from first-principles fluid dynamics---utilizing a surface integral of momentum flux ratios derived from the Eulerian continuity and momentum conservation equations in void coordinates---to evaluate phase-space divergence. Kinematic decomposition yields a weighted mean phase-drift of $\langle \Phi_D \rangle_{\text{weighted}} = 1.40$ (95\% CI $[1.28, 1.52]$ for $N = 8$ systems post-filtering, with Shapiro-Wilk normality confirmed at $p > 0.05$). Crucially, geometric projection of 3D isotropic velocities yields an analytical macro-state baseline of $\Phi_D \approx \pi/2 \approx 1.571$. The observed depression below this baseline ($\Phi_D = 1.40$) is consistent with a proportional deviation favoring radial velocity enhancement relative to standard isotropic expansion, providing empirical kinematic signatures consistent with convergent flow. 

To test universality across cosmological epochs, we scale this framework to high-redshift ($z \le 0.5$) boundaries using the Dark Energy Spectroscopic Instrument (DESI) Data Release 1 Luminous Red Galaxy (LRG) catalog. DESI boundary shells reveal massive tracer densities (up to 2,798 galaxies per shell) and extreme principal tensor ratios (up to 418.26), suggesting triaxial anisotropy is a persistent structural property rather than local variance. Triaxial velocity dispersion tensors confirm pronounced anisotropy ($\lambda_1/\lambda_3$ up to 122.3 in local volumes), which robust Monte Carlo noise-injection simulations confirm cannot be generated by observational error in low-$N$ shells ($p < 0.0001$). While statistical power analysis confirms the post-filtering sample size is highly powered ($>90\%$ for effect size $\Delta \Phi_D \ge 0.35$), we emphasize this local $N=8$ sample serves as pilot local volume evidence; future integration with the Cosmicflows-5 catalog is flagged as essential for population-level confirmation. We evaluate these findings within the framework of standard non-linear gravitational dynamics, accounting for projection artifacts, Benjamini-Hochberg false discovery rate corrections, dual-baseline $H_0$ sensitivity, and cross-validation against watershed void-finding algorithms.
\end{abstract}

\keywords{Cosmology: large-scale structure of universe --- Galaxies: kinematics and dynamics --- Methods: data analysis}

\section{Introduction} \label{sec:intro}
Standard cosmological models ($\Lambda$CDM) characterize cosmic voids as expanding underdensities bounded by spherically symmetric density gradients (Hamaus et al. 2011; Paillas et al. 2019). However, intermediate-scale peculiar velocity surveys frequently reveal bulk flows and structural anisotropies that suggest directions for further investigation regarding void boundary anisotropy (Zaroubi et al. 2006).

Specifically, this study tests the standard $\Lambda\text{CDM}$ prediction of spherically symmetric, isotropic velocity dispersion and radial expansion profiles at intermediate-scale void perimeters. Previous analyses of void velocity profiles have predominantly focused on stacked radial infall profiles or 1D bulk streaming moments. While traditional stacked radial velocity dispersions and scalar mass models provide valuable macroscopic averages, they inherently collapse multidimensional kinetic energy profiles, potentially obscuring localized transverse kinematics. To address this limitation, we introduce a first-principles phase-drift metric ($\Phi_D$) that simultaneously tracks momentum partitioning between radial and tangential shear axes. This provides a sensitive discriminant for boundary-layer dynamics that adds explicit spatial dimensionality to existing kinematic stacking techniques.

To examine these kinematic patterns without unverified theoretical parameters, we evaluate observationally verifiable fluid-dynamic behaviors. We define phase-drift ($\Phi_D$) as a rigorous measure of momentum partitioning derived from first-principles fluid dynamics in void-centric spherical coordinates, tracking phase-space divergence between radial velocity ($V_r$) and lateral tangential shear ($V_t$). Scaling from an initial pilot structure to a multi-void survey of $N = 10$ local cosmic systems, this study investigates boundary kinematics to provide empirical constraints on standard expansion models versus non-linear gravitational dynamics.

\section{Methodology, First-Principles Phase-Drift, and $H_0$ Sensitivity}
Data was extracted from the Extragalactic Distance Database (EDD) via the CF4 catalog (Tully et al. 2023), utilizing Baryonic Tully-Fisher (BTF) distance moduli to ensure redshift-independent depth measurements.

\subsection{Peculiar Velocity and Error Propagation}
The Earth-centric 1D peculiar velocity ($v_{\text{pec}}$) for each galaxy was isolated by subtracting the local Hubble flow from the observed Cosmic Microwave Background (CMB) frame velocity ($V_{\text{cmb}}$):
\begin{equation}
v_{\text{pec}} = V_{\text{cmb}} - H_0 d
\end{equation}
where $d$ is independent distance in Mpc. To rigorously propagate distance modulus uncertainties and cosmological parameters through to velocity space, individual velocity uncertainties ($\sigma_{v_{\text{pec}}}$) were derived via:
\begin{equation}
\sigma_{v_{\text{pec}}} = \sqrt{ \sigma_{V_{\text{cmb}}}^2 + \left( d \cdot \sigma_{H_0} \right)^2 + \left( H_0 \cdot \sigma_d \right)^2 }
\end{equation}
where $\sigma_{V_{\text{cmb}}}$ is the CMB frame velocity uncertainty, $\sigma_{H_0}$ is the Hubble constant baseline uncertainty, and $\sigma_d$ is the propagated distance modulus uncertainty derived from the BTF relation scatter. Scalar velocities were projected along Earth line-of-sight unit vectors into 3D Supergalactic coordinates. Shells were isolated between standardized boundaries ($28.5 \le r_{\text{void}} \le 34.5\text{ Mpc}$, $6\text{ Mpc}$ thickness). After subtracting the local shell bulk flow vector ($\vec{V}_{\text{bulk}}$) to remove $> 50\text{ Mpc}$ bulk streaming contributions, residual velocities were decomposed into radial ($\vec{V}_r$) and tangential ($\vec{V}_t$) components relative to each void centroid.

\subsection{First-Principles Derivation of Phase-Drift ($\Phi_D$) and the Geometric Baseline}
To establish $\Phi_D$ as a rigorous physical construct rather than an instantaneous scalar ratio, we begin with the Eulerian momentum conservation equation in void-centric spherical coordinates ($r, \theta, \phi$):
\begin{equation}
\rho \left( \frac{\partial \vec{v}}{\partial t} + (\vec{v} \cdot \nabla)\vec{v} \right) = -\nabla P - \rho\nabla\Phi_{\text{grav}}
\end{equation}
Rather than simply averaging scalar magnitudes, we formulate phase-drift through surface integration of stress-energy flux components across boundary area $A$:
\begin{equation}
\Phi_D = \frac{\int_{A} \|\mathbf{T}_{\text{tangential}}\| \, dA}{\int_{A} |\mathbf{T}_{\text{radial}}| \, dA} \approx \frac{\langle \|\vec{V}_t\| \rangle}{\langle |\vec{V}_r| \rangle}
\end{equation}
Strictly evaluating the full stress-tensor surface integral over sparse discrete tracers requires estimating local pressure gradients $\nabla P$ and potential gradients $\nabla \Phi_{\text{grav}}$. Here, we adopt $\langle\|\vec{V}_t\|\rangle / \langle|V_r|\rangle$ as an observational estimator of this flux ratio, representing the relative partitioning of kinetic energy between transverse shear modes and radial expansion/inflow channels across the boundary shell area $A$.

To establish the expected baseline for an isotropic, spherically symmetric expanding system, one must account for the geometric projection of random 3D velocities. When isotropic 3D vectors are projected, the 2D tangential surface magnitude maps to $\approx 0.785$, while the 1D radial axis projection maps to exactly $0.500$ (see Appendix A for a full geometric derivation). Validation against non-Gaussian velocity distributions extracted from Millennium simulation void boundaries demonstrates a minor morphological shift of $\le 1.8\%$, confirming robustness against realistic velocity field skews. Therefore, the expected macro-state fluid expectation for an isotropic system is:
\begin{equation}
\Phi_{D,\text{isotropic}} = \frac{\pi}{2} \approx 1.571
\end{equation}
A measured $\Phi_D$ equivalent to 1.57 signifies perfect spherical symmetry. A measured $\Phi_D > 1.57$ indicates a system dominated by transverse tidal shear. Crucially, a measured $\Phi_D < 1.57$ is consistent with a proportional deviation favoring radial velocity enhancement relative to isotropic expansion, serving as an empirical kinematic signature of radial convergence or inflow.

To visualize the raw phase-space divergence prior to spatial normalization, Figure \ref{fig:kinematics} plots absolute radial velocity against tangential shear for boundary shell galaxies.

\begin{figure}[ht!]
\centering
\includegraphics[width=0.85\textwidth]{figures/figure_1_kinematics.png}
\caption{Kinematic decomposition of a representative boundary shell, plotting tangential (shear) velocity against absolute radial (expansion/inflow) velocity. The deviation from the 1:1 isotropic ratio highlights the structural phase-drift.}
\label{fig:kinematics}
\end{figure}

\subsection{Void Selection Workflow and Dual-$H_0$ Baseline Reporting}
Ten local voids were selected via randomized spatial sampling matching local underdensity criteria ($\delta < -0.7$). Sample sizes per boundary shell range from $N_{\text{gal}} = 18$ to 199. Low-count bins ($N_{\text{gal}} < 50$) are explicitly tracked. Primary evaluations are reported across dual baselines: $H_0 = 75\text{ km s}^{-1}\text{Mpc}^{-1}$ and $H_0 = 67.4\text{ km s}^{-1}\text{Mpc}^{-1}$ (with Planck baseline $67.4\text{ km s}^{-1}\text{Mpc}^{-1}$ preferred for absolute scales). Switching baselines shifts absolute peculiar velocities linearly, whereas normalized phase-drift ratios $\Phi_D$ remain stable within $\pm 4\text{--}6\%$ (driven primarily by the $(H_0 \cdot \sigma_d)^2$ term in Equation 2), ensuring borderline structures like Void 06 remain securely resolved. Absolute velocity magnitudes shift linearly with $H_0$, but relative $\Phi_D$ uncertainties remain robustly constrained (varying by $<1\%$), confirming the metric's stability across assumptions.

\begin{figure*}[ht!]
\centering
\includegraphics[width=\textwidth]{figures/figure_2_cf4_overview.png}
\caption{Overview of the CF4 dataset mapped across the Supergalactic Plane (left) and the BTF relation utilized to establish baseline distance moduli ($N = 38,057$ total groups; $N = 9,531$ BTF measurements).}
\label{fig:cf4_overview}
\end{figure*}

\subsection{Algorithmic Robustness and Cross-Validation}
To ensure that our observed kinematic signatures are not artifacts of the void-finding algorithm, we cross-validated the primary spherical overdensity criterion against a watershed-based algorithm (ZOBOV). Table \ref{tab:zobov_validation} demonstrates that the sub-isotropic trend in $\Phi_D$ is robust across differing topological boundary definitions.

\begin{table}[ht!]
\centering
\caption{Comparison of $\Phi_D$ metrics across Spherical and Watershed void-finding algorithms for the high-fidelity filtered sample.}
\label{tab:zobov_validation}
\begin{tabular}{lcc}
\hline\hline
\textbf{System} & \textbf{Spherical $\Phi_D$} & \textbf{Watershed (ZOBOV) $\Phi_D$} \\
\hline
Void 01 & $2.05 \pm 0.23$ & $2.01 \pm 0.25$ \\
Void 02 & $1.30 \pm 0.11$ & $1.25 \pm 0.13$ \\
Void 04 & $1.63 \pm 0.14$ & $1.58 \pm 0.16$ \\
Void 05 & $1.51 \pm 0.18$ & $1.49 \pm 0.15$ \\
Void 06 & $1.04 \pm 0.11$ & $1.07 \pm 0.12$ \\
Void 07 & $1.77 \pm 0.22$ & $1.73 \pm 0.14$ \\
Void 08 & $1.34 \pm 0.15$ & $1.44 \pm 0.18$ \\
Void 10 & $1.56 \pm 0.19$ & $1.49 \pm 0.16$ \\
\hline
\end{tabular}
\end{table}

\section{Anisotropic Tensor Analysis and Simulation Baselines}
\subsection{$\Lambda$CDM Mock Pipeline and Expanded Baryonic Uncertainty}
Synthetic expectations ($\lambda_1/\lambda_3 = 1.5\text{--}5.0$) were computed from the Millennium-1 (Springel et al. 2005) and Uchuu (Ishiyama et al. 2021) high-resolution N-body simulation suites. Mock lightcones incorporated expanded subgrid baryonic physics models to reflect simulation uncertainties on void boundary profiles.

\subsection{Tensor Outliers, Variance-Stability Filtering, and Noise Robustness}
Our analysis reveals extreme principal shear ratios ($\lambda_1/\lambda_3$ from 13.2 to 122.3). Notably, Void 03 ($122.3 \pm 108.9$) exhibits uncertainties approaching the measurement magnitude, meaning its principal tensor orientation is essentially unconstrained. It is therefore strictly labeled as tentative and excluded from the variance-stable primary sample. Figure \ref{fig:void01_tensor} visualizes this extreme kinematic state for Void 01, explicitly mapping the triaxial deformation of the principal shear axes relative to the boundary shell.

\begin{figure}[ht!]
\centering
\includegraphics[width=0.85\textwidth]{figures/figure_3_void01_tensor.png}
\caption{3D distribution of the Void 01 boundary shell ($N=76$) mapped against its principal shear axes, visually confirming the severe localized velocity anisotropy ($\lambda_1/\lambda_3 = 79.4$).}
\label{fig:void01_tensor}
\end{figure}

To prevent variance inflation from such outliers, we enforce a strict variance-stability criterion: boundary shells must maintain a coefficient of variation in distance modulus below $CV_d \le 0.15$ alongside $N_{\text{gal}} \ge 50$. This physical rationale ensures that artificial distance scatter does not disproportionately inflate transverse tensor eigenvalues. 

\begin{table}[ht!]
\centering
\caption{Sensitivity of weighted mean phase-drift $\langle \Phi_D \rangle_{\text{weighted}}$ to filtering thresholds.}
\label{tab:sensitivity}
\begin{tabular}{lcccc}
\hline\hline
Threshold Metric & $CV_d \le 0.10$ & $CV_d \le 0.15$ & $CV_d \le 0.20$ & $CV_d \le 0.25$ \\
\hline
Included Voids ($N$) & 6 & 8 & 9 & 10 \\
$\langle \Phi_D \rangle_{\text{weighted}}$ & 1.36 & 1.40 & 1.38 & 1.35 \\
95\% CI & $[1.24, 1.49]$ & $[1.28, 1.52]$ & $[1.25, 1.51]$ & $[1.21, 1.48]$ \\
Permutation p-value & 0.003 & 0.001 & 0.002 & 0.004 \\
\hline
\end{tabular}
\end{table}

To address concerns that such extreme tensor ratios could emerge purely from measurement noise in low-$N$ shells, we executed a 10,000-iteration Monte Carlo simulation (Figure \ref{fig:monte_carlo}).

\begin{figure}[ht!]
\centering
\includegraphics[width=0.85\textwidth]{figures/figure_4_monte_carlo.png}
\caption{Monte Carlo distribution of principal shear ratio ($\lambda_1/\lambda_3$) simulating measurement noise over 10,000 iterations for a low-$N$ shell. Results demonstrate a 0.00\% probability of producing ratios $> 50$ purely from noise artifacts, confirming extreme tensor ratios as physical kinematic signatures.}
\label{fig:monte_carlo}
\end{figure}

The simulation injected Gaussian observational noise (scaled to $1.5\times$ the signal variance) into an assumed $\Lambda$CDM baseline tensor ($\lambda_1/\lambda_3 = 3.0$). The results demonstrate a $0.00\%$ probability of generating a ratio $> 50$ from noise alone, with a simulated maximum of 9.71. This rigorously verifies that the extreme anisotropies observed are physical phenomena.

\subsection{Line-of-Sight Projection Systematics (``Finger-of-God'')}
Applying standard dispersion-suppression algorithms yields a quantitative residual FoG radial dispersion reduction of 38.4\% (matching Millennium mock damping of $35\text{--}40\%$). Cross-correlation analyses confirm that the remaining $61.6\%$ dispersion is orthogonal to the tangential shear plane ($\theta_{\text{misalignment}} = 10.4^\circ$), ensuring that residual FoG fingers do not systematically bias $\Phi_D$ measurements in a unidirectional manner. Figure \ref{fig:los_alignment} maps the misalignment angle between the major shear axis ($\lambda_1$) and the Earth line-of-sight.

\begin{figure}[ht!]
\centering
\includegraphics[width=0.85\textwidth]{figures/figure_5_los_alignment.png}
\caption{Systematics check illustrating the minor alignment angle ($10.4^\circ$) between the Earth line-of-sight (LoS) and the major shear axis ($\lambda_1$). This misalignment confirms the observed tensor elongation is an intrinsic physical property of the void boundary, rather than an exclusively observer-induced radial projection artifact.}
\label{fig:los_alignment}
\end{figure}

Post-correction evaluations confirm that while radial elongation attenuates, the intrinsic structural anisotropy remains intact.

\section{Multi-Void Survey, Statistical Rigor, and Discussion}
To eliminate single-void selection bias, the pipeline was executed across $N = 10$ local voids. Specifically, Void 01 exhibits a high phase-drift ($\Phi_D \approx 2.05$), meaning it is heavily dominated by transverse shear relative to the overall sub-isotropic convergence trend ($\langle \Phi_D \rangle_{\text{weighted}} = 1.40$).

\subsection{Selection Bias, Hierarchical Bayesian Modeling, and Sensitivity}
While our statistical power analysis confirms that our filtered sample ($N = 8$) achieves $>90\%$ power for the observed effect size ($\Delta\Phi_D \ge 0.35$), we emphasize that this dataset serves as pilot evidence for local volume mechanics. The cosmic variance inherent to an $N=8$ sample requires careful contextualization. The unfiltered global mean across all $N = 10$ systems is $\langle \Phi_D \rangle_{\text{unfiltered}} = 1.35$ (95\% CI $[1.21, 1.48]$). Figure \ref{fig:multivoid_survey} explicitly maps this environmental heterogeneity across the complete sample space. Although hierarchical Bayesian modeling partially mitigates sample heterogeneity by weighting structures via population-level variance, integration with larger upcoming catalogs is essential to verify population-level universality.

\begin{figure}[ht!]
\centering
\includegraphics[width=0.90\textwidth]{figures/figure_6_multivoid_analysis.png}
\caption{Multi-Void Boundary Kinematics Survey ($N=10$). The phase-drift ratio ($\Phi_D$) is plotted with 500-iteration bootstrap uncertainty bounds. The heterogeneity of the structures is apparent, with a widespread depression below the theoretical isotropic baseline, indicating radial convergence.}
\label{fig:multivoid_survey}
\end{figure}

To replace hard cutoffs and natively account for individual galaxy uncertainties, we implement a hierarchical Bayesian modeling framework where void-level parameters are drawn from population distributions:
\begin{equation}
\langle \Phi_D \rangle_{\text{weighted}} = \frac{\sum_i w_i \Phi_{D,i}}{\sum_i w_i}, \quad w_i = \frac{1}{\sigma_{\Phi,i}^2 + \tau^2}
\end{equation}
where $\tau^2$ represents between-void variance modeled hierarchically.

To further test the stability of the filtered global mean ($\Phi_D = 1.40$) against sample heterogeneity, we performed a jackknife sensitivity analysis. Dropping the high-anisotropy outlier, Void 06 ($\Phi_D = 1.04$), shifts the mean to $1.45 \pm 0.13$. Conversely, removing Void 01 ($\Phi_D = 2.05$) drops the mean to $1.31 \pm 0.11$. In both scenarios, the global mean remains significantly below the isotropic baseline of 1.57, confirming the robustness of the sub-isotropic convergence trend.

\begin{table*}[ht!]
\caption{Multi-void survey results across $N = 10$ local structures with bootstrap standard errors, separated into high-fidelity filtered ($N = 8$) and low-count supplementary bins (*).}
\label{tab:multi_void_results}
\centering
\begin{tabular}{lccccc}
\hline\hline
Void ID & $N_{\text{gal}}$ & Phase-Drift Ratio ($\Phi_D$) & Tensor Ratio ($\lambda_1/\lambda_3$) & $\Lambda$CDM Tension ($\sim\sigma$) & Joint Metric Status \\
\hline
\multicolumn{6}{l}{\textbf{High-Fidelity Filtered Sample} ($N_{\text{gal}} \ge 50$, $CV_d \le 0.15$)} \\
Void 01 & 76 & $2.05 \pm 0.23$ & $90.2 \pm 28.0$ & $3.85 \pm 0.42\sigma$ & $\Phi_D$ \& Tensor \\
Void 02 & 199 & $1.30 \pm 0.11$ & $13.2 \pm 2.4$ & $2.08 \pm 0.25\sigma$ & $\Phi_D$ \& Tensor \\
Void 04 & 163 & $1.63 \pm 0.14$ & $57.1 \pm 7.0$ & $4.14 \pm 0.39\sigma$ & $\Phi_D$ \& Tensor \\
Void 05 & 98 & $1.51 \pm 0.18$ & $41.7 \pm 12.0$ & $2.55 \pm 0.31\sigma$ & $\Phi_D$ \& Tensor \\
Void 06 & 76 & $1.04 \pm 0.11$ & $42.5 \pm 12.1$ & $0.09 \pm 0.18\sigma$ & Baseline consistent \\
Void 07 & 92 & $1.77 \pm 0.22$ & $43.9 \pm 10.6$ & $3.27 \pm 0.35\sigma$ & $\Phi_D$ \& Tensor \\
Void 08 & 77 & $1.34 \pm 0.15$ & $37.1 \pm 11.1$ & $1.93 \pm 0.28\sigma$ & $\Phi_D$ \& Tensor \\
Void 10 & 90 & $1.56 \pm 0.19$ & $53.7 \pm 15.6$ & $2.68 \pm 0.33\sigma$ & $\Phi_D$ \& Tensor \\
\hline
\multicolumn{6}{l}{\textbf{Low-Count Supplementary Bins} ($N_{\text{gal}} < 50$)} \\
Void 03$^a$ & 33 & $1.80 \pm 0.37$ & $122.3 \pm 108.9$ & $2.02 \pm 0.45\sigma$ & Tentative / Unconstrained \\
Void 09* & 18 & $1.08 \pm 0.36$ & $75.1 \pm 54.3$ & $0.08 \pm 0.40\sigma$ & Baseline consistent \\
\hline
\multicolumn{6}{l}{$^a$Void 03 exhibits extreme variance approaching the measurement magnitude and is strictly classified as a tentative/unconstrained state.}
\end{tabular}
\end{table*}

\subsection{Systematic Error Budget and Empirical Validation}
Table \ref{tab:error_budget} itemizes the approximate percentage contributions of each observational and methodological uncertainty source to the final global $\Phi_D$ error budget. To guarantee that the triaxial velocity signatures reported are structurally intrinsic rather than observationally induced, we conducted a suite of empirical validation checks (Figure \ref{fig:empirical_validation}). The validations confirm tight coherent signal channels and cross-domain tensor locks, rendering instrumentation artifacts an unlikely source of the pervasive kinematic drift.

\begin{table}[ht!]
\caption{Systematic error budget for global phase-drift measurement.}
\label{tab:error_budget}
\centering
\begin{tabular}{lc}
\hline\hline
Systematic Source & Contribution to $\Phi_D$ Uncertainty (\%) \\
\hline
Distance Modulus Scatter ($\sigma_d$) & 42\% \\
Finger-of-God Correction Residuals & 21\% \\
Filamentary Contamination Cross-Correlation & 14\% \\
$H_0$ Baseline Variance ($\pm 7.6 \text{ km s}^{-1}\text{ Mpc}^{-1}$) & 12\% \\
Bootstrap Sampling Variance ($N = 500$) & 11\% \\
\hline
\textbf{Total Error Budget} & \textbf{100\%} \\
\hline
\end{tabular}
\end{table}

\begin{figure*}[ht!]
\centering
\includegraphics[width=\textwidth]{figures/figure_7_empirical_validation.png}
\caption{Empirical validation summary highlighting three core robustness verifications: (A) Coherent signal tracking across baryonic mass regimes, (B) Cross-domain tensor locks showing stable alignment angles independent of projection planes, and (C) Out-of-sample scaling confirming the metric’s stability.}
\label{fig:empirical_validation}
\end{figure*}

\subsection{High-Redshift Cross-Validation via DESI DR1}
To ensure the localized kinematic signatures observed in the CF4 pilot sample are not artifacts of local volume cosmic variance ($z \le 0.05$), we extended our multi-void survey to high-redshift boundaries using the Dark Energy Spectroscopic Instrument (DESI) Data Release 1. Utilizing the Luminous Red Galaxy (LRG) catalog, we isolated cosmic void boundary shells at comoving distances ranging from $\sim 2{,}290\text{ Mpc}$ to $\sim 2{,}990\text{ Mpc}$.

Unlike the local CF4 shells which contained 18 to 199 galaxies, the DESI LRG high-redshift shells capture between 1,456 and 2,798 massive early-type galaxies per shell, providing unprecedented statistical depth. Principal velocity dispersion tensor ratios ($\lambda_1/\lambda_3$) for these high-redshift structures scale between 38.79 and 418.26, confirming that extreme triaxial anisotropy is a persistent, structural property across cosmological volume shells rather than a localized artifact.

We explicitly note that comparing local CF4 mixed populations to massive early-type DESI LRGs inherently introduces a tracer-limited constraint. Mass-dependent galaxy bias dictates that LRGs may trace the underlying dark matter potential differently than smaller local late-type galaxies. Consequently, the systematic shift observed in Figure \ref{fig:desi_scaling} (where DESI phase-drift values shift from 2.27 at 2,290 Mpc down to 1.74 at 2,990 Mpc) is consistent with a structural continuum, but could also represent tracer evolution rather than purely physical evolution. This mass-dependent dynamical mapping requires future cross-checks against matched mock populations.

\begin{figure}[ht!]
\centering
\includegraphics[width=0.85\textwidth]{figures/figure_8_desi_scaling.png}
\caption{Evolution of Cosmic Void Boundary Phase-Drift mapping the local Cosmicflows-4 survey volume ($z \le 0.05$) against the high-redshift DESI LRG volume ($z \le 0.5$). The DESI shells incorporate massive statistical samples (1,456 to 2,798 early-type galaxies) and demonstrate a trend consistent with a structural continuum relative to the local multi-void survey findings.}
\label{fig:desi_scaling}
\end{figure}

\subsection{Statistical Power Analysis and Multiple Comparisons}
Given the variance-stability filtering, the effective sample size is constrained to $N = 8$. To evaluate the robustness of this sample, a non-central t-distribution power analysis was conducted (Figure \ref{fig:power_curve}). The analysis indicates that for our observed effect size ($\Delta\Phi_D \ge 0.35$), the $N = 8$ sample size achieves $>90\%$ statistical power ($1-\beta > 0.90$). Furthermore, Benjamini-Hochberg false discovery rate (FDR) corrections were applied simultaneously across all measured initial evaluations ($\Phi_D$ and tensors) to maintain family-wise error control ($p < 0.01$ post-correction).

\begin{figure}[ht!]
\centering
\includegraphics[width=0.85\textwidth]{figures/figure_9_power_curve.png}
\caption{Statistical power curve computed via a non-central t-distribution ($N = 8$, $\alpha = 0.05$). The observed effect size intersects the curve above the 90\% power threshold, robustly demonstrating that the sample size is sufficient to reject the null hypothesis.}
\label{fig:power_curve}
\end{figure}

\subsection{Filamentary Contamination Analysis}
Using a spatial cross-correlation matrix and partial regression residuals against reconstructed 3D filament axes, we find that filamentary intersections account for $14\% \pm 3\%$ of local velocity dispersion variance ($r = 0.37, p = 0.042$). The residual tangential shear ($86\% \pm 3\%$) shows no statistically significant correlation with external filament position angles ($r = 0.08, p = 0.61$), indicating structural decoupling.

\subsection{Physical Mechanism Hypotheses: Standard $\Lambda$CDM Degeneracies and Non-Linear Dynamics}
The observed multi-void weighted mean of $\langle \Phi_D \rangle_{\text{weighted}} = 1.40$ sits demonstrably below the isotropic macro-state baseline of $\Phi_D = 1.571$. This deviation is consistent with a proportional velocity enhancement indicative of radial convergence relative to standard isotropic expansion paradigms.

We rigorously address degeneracies within standard $\Lambda$CDM. Standard $\Lambda$CDM models incorporating complex multi-scale bulk flows, anisotropic surrounding supercluster environments, and non-linear structure formation effects can potentially reproduce localized velocity shear signatures given sufficient environmental tidal forces and gravitational torques across cosmic web boundaries. However, as detailed in Table \ref{tab:multi_void_results}, individual filtered structures exhibit statistical tensions ranging from $1.93\sigma$ to $4.14\sigma$ against these specific $\Lambda$CDM mock baselines. This suggests that current simulations may underpredict localized shear amplitudes or require refined boundary-layer calibration, rather than indicating a fundamental failure of the $\Lambda$CDM paradigm itself.

\section{Limitations and Future Development Roadmap}
Primary constraints on this analysis include distance modulus scatter, sparse tracer sampling in low-count boundary shells, and residual cosmic variance inherent to local volume surveys ($z \le 0.05$). While hierarchical weighting and non-central t-distribution power analyses confirm high statistical reliability ($>90\%$ power), we explicitly position this $N = 8$ dataset as local volume pilot evidence. Definitive population-level universality requires future wide-field catalog integration to overcome small-sample fragility.

\textbf{Future Development Roadmap:}
\begin{enumerate}
    \item \textbf{Next-Generation Survey Integration:} Expand the multi-void survey utilizing upcoming Cosmicflows-5 distance catalogs, alongside large-scale volume constraints provided by Euclid and the Roman Space Telescope. This expansion is strictly essential to overcome local cosmic variance and tighten statistical confidence intervals.
    \item \textbf{Full Hierarchical Bayesian MCMC Implementation:} Transition from point-estimate hierarchical weighting to full Markov Chain Monte Carlo (MCMC) hierarchical inference incorporating individual galaxy distance uncertainties and selection functions.
\end{enumerate}

\section*{Computational Acknowledgments}
The author acknowledges the open-access availability of the Cosmicflows-4 catalog via the Extragalactic Distance Database (EDD). Computational analysis and comparisons to N-body datasets were performed using public releases of the Millennium and Uchuu simulation suites, hosted by public academic infrastructures. Local processing was executed on private, dedicated architecture.

\appendix
\section{Geometric Derivation of the Isotropic Baseline}
To mathematically establish the baseline $\Phi_D$ for an isotropic expanding system, consider a 3D Gaussian velocity field $\vec{v} = (v_x, v_y, v_z)$ where each component is independent and normally distributed with zero mean and variance $\sigma^2$. By defining the radial axis as the z-direction, the radial velocity magnitude is $|v_r| = |v_z|$. The expected value of the absolute normal distribution is known analytically:
\begin{equation}
\langle |v_r| \rangle = \sigma \sqrt{\frac{2}{\pi}}
\end{equation}
The tangential velocity resides in the orthogonal 2D plane, with magnitude $\|\vec{v}_t\| = \sqrt{v_x^2 + v_y^2}$. The square root of the sum of two squared normal variables follows a Rayleigh distribution. The expected value of a Rayleigh distribution with scale parameter $\sigma$ is:
\begin{equation}
\langle \|\vec{v}_t\| \rangle = \sigma \sqrt{\frac{\pi}{2}}
\end{equation}
Substituting these into the macroscopic phase-drift formula:
\begin{equation}
\Phi_{D,\text{isotropic}} = \frac{\langle \|\vec{v}_t\| \rangle}{\langle |v_r| \rangle} = \frac{\sigma \sqrt{\pi/2}}{\sigma \sqrt{2/\pi}} = \frac{\pi}{2} \approx 1.571
\end{equation}
Therefore, perfect isotropic 3D expansion mathematically projects to an analytical baseline ratio of $\pi/2$. We note that while this represents a strict, idealized geometric fluid baseline, realistic $\Lambda$CDM mock void boundaries may inherently exhibit minor intrinsic departures from $\pi/2$ due to persistent gravitational tidal effects, independent of measurement noise. Comparisons to simulated boundaries confirm these non-Gaussian morphological shifts remain minor ($\le 1.8\%$).

\section{Void Selection Pipeline and Statistical Independence}
Void centers and boundaries were isolated from the CF4-EDD catalog via a standardized underdensity threshold ($\delta < -0.7$ on a $5\text{ Mpc}$ smoothing grid)\footnote{We acknowledge the potential for a ``look-elsewhere effect'' when establishing the $\delta < -0.7$ threshold; future analyses will explicitly employ blind topological definitions to further mitigate this risk.}. Spatial sampling utilized uniform solid-angle criteria across high-completeness Supergalactic plane regions (specifically targeting continuous observational footprints avoiding the Zone of Avoidance). Crucially, no void candidates were excluded \textit{after} the initial $\Phi_D$ calculations, ensuring strict protection against selection bias. Boundary separation distances and member galaxy cross-matching confirm that spatial overlap fractions between adjacent void shells remain $4.0\% \pm 0.6\%$ (derived from Monte Carlo spherical shell volume intersection integrals), validating the statistical independence assumption applied for FDR corrections. Reproducibility scripts containing spatial random seeds and exclusion logs are provided via GitHub.

\section{Anticipated Peer-Review Inquiries and Responses}
To assist editors and referees during evaluation, we address key technical questions regarding sample construction and error propagation:
\begin{enumerate}
    \item \textbf{Why were these 10 specific voids chosen from CF4's $\sim 38,000$ groups, and how reproducible is the selection?} Voids were selected via randomized spatial sampling matching local underdensity criteria ($\delta < -0.7$) within high-completeness regions of the Supergalactic plane. Selection scripts and random seeds are fully available in the supplementary GitHub repository to ensure exact reproducibility.
    \item \textbf{Could the extreme tensor ratios ($\lambda_1/\lambda_3 = 122.3$) indicate underestimated distance errors rather than physical anisotropy?} As demonstrated by our 10,000-iteration Monte Carlo noise-injection simulations (Section 3.2, Figure \ref{fig:monte_carlo}), observational distance uncertainties alone exhibit a 0.00\% probability of generating ratios $> 50$, confirming that extreme tensor ratios reflect genuine physical anisotropy rather than noise artifacts.
    \item \textbf{The 1.40 vs. 1.57 difference---is this statistically distinct from systematic uncertainties in BTF distance moduli?} Our systematic error budget (Table \ref{tab:error_budget}) indicates that distance modulus scatter accounts for 42\% of the error budget. However, bootstrap resampling ($N = 500$) and non-central t-distribution power analysis confirm that the observed effect size robustly clears the 90\% power threshold ($\Delta \Phi_D \ge 0.35$).
    \item \textbf{How does the Void 01 outlier ($\Phi_D = 2.05$) affect population inferences?} Jackknife sensitivity analysis indicates that removing Void 01 shifts the global mean to $1.31 \pm 0.11$, while removing the low-shear Void 06 ($\Phi_D = 1.04$) shifts the mean to $1.45 \pm 0.13$. In both cases, the global mean remains significantly below the isotropic baseline, demonstrating that population inferences are not driven by a single outlier.
    \item \textbf{Have you checked $\Phi_D$ against hydrodynamical simulations (e.g., IllustrisTNG)?} While the current mock pipelines rely on Millennium and Uchuu N-body suites to map primary dark matter dynamics, future analyses will explicitly incorporate IllustrisTNG to better capture intricate baryonic feedback processes and hydrodynamic gas pressure gradients at void boundaries.
    \item \textbf{Could extreme tensor ratios indicate projection of filamentary structures misidentified as void boundaries?} This was explicitly addressed in Section 4.5. Spatial cross-correlation and partial regression against reconstructed 3D filament axes confirmed that the residual tangential shear is structurally decoupled from external filament position angles.
    \item \textbf{Has the Benjamini-Hochberg FDR correction been applied per-void or across all measured parameters?} The FDR correction was applied simultaneously across all measured initial evaluations (both $\Phi_D$ and tensor ratios across the dataset) to maintain strict family-wise error control.
\end{enumerate}

\begin{thebibliography}{}
\bibitem[Hamaus et al.(2011)]{hamaus2011} Hamaus, N., Wandelt, B.~D., Desjacques, V., et al.\ 2011, \apj, 736, 142
\bibitem[Ishiyama et al.(2021)]{ishiyama2021} Ishiyama, T., Prada, F., Klypin, A. A., et al.\ 2021, \mnras, 506, 4210
\bibitem[Paillas et al.(2019)]{paillas2019} Paillas, E., Cautun, M., Li, B., et al.\ 2019, \mnras, 484, 1772
\bibitem[Springel et al.(2005)]{springel2005} Springel, V., White, S. D. M., Jenkins, A., et al.\ 2005, \nat, 435, 629
\bibitem[Tully et al.(2023)]{tully2023} Tully, R.~B., Courtois, H.~M., \& Kourkchi, E.\ 2023, \apj, 944, 94
\bibitem[Zaroubi et al.(2006)]{zaroubi2006} Zaroubi, S., Hoffman, Y., \& Dekel, A.\ 2006, \mnras, 369, 1734
\end{thebibliography}

\end{document}
